% TOP_PROBLEM: {"id":30007551,"problem_number":"LOCAL-30007551","title":"Redirecting technological change","category_id":21,"category":{"id":21,"name":"economics","display_name":"Economics","description":"Open economic theory statements, published research agendas, and proposed scoped specifications; question provenance and historical priority are qualified per record.","slug":"economics","order_index":21,"created_at":"2026-10-08T00:00:00Z"},"status":"research_question","external_url":null,"legacy_ids":[],"published":false,"statement":"Design innovation-direction incentives that maximize social surplus under private information, fiscal limits, capture risk, and experimentation requirements.","economics":{"original_id":40,"field":"Growth and productivity","kind":"design","formalization_basis":"adapted_research_question","source_status":"explicit_open","priority_status":"earliest_found","priority_note":"This is the earliest explicit relevant statement verified in this audit, not a claim of original priority; older literature and earlier versions may contain antecedents.","earliest_verified_reference":{"authors":["Ufuk Akcigit"],"title":"Creative destruction in economic growth","year":2026,"url":"https://onlinelibrary.wiley.com/doi/10.1111/sjoe.70037","doi":"10.1111/sjoe.70037","locator":"Section 7, “Open questions,” item (7)","evidence_summary":"The author explicitly asks how policy can redirect innovation toward social priorities when private and social returns differ while preserving experimentation."},"current_reference":{"authors":["Ufuk Akcigit"],"title":"Creative destruction in economic growth","year":2026,"url":"https://onlinelibrary.wiley.com/doi/10.1111/sjoe.70037","doi":"10.1111/sjoe.70037","locator":"Section 7, “Open questions,” item (7)","evidence_summary":"The author explicitly asks how policy can redirect innovation toward social priorities when private and social returns differ while preserving experimentation."},"formalization_note":"The source explicitly asks about directing innovation toward social needs while preserving experimentation. The two-direction mechanism and capture constraints are adapted."},"statusQualification":"Published question or research agenda verified at the cited locator. The mathematical specification is adapted for this catalogue and includes additional assumptions; source publication does not establish first-ever priority or certify this exact model remains unsolved. The source explicitly asks about directing innovation toward social needs while preserving experimentation. The two-direction mechanism and capture constraints are adapted.","statusReviewedAt":"2026-10-08"}
% FIRST_POSTED: 2026-10-08
% ENTRY_KIND: statement_only
% !TeX program = lualatex
\documentclass[11pt]{article}
\usepackage{fontspec}
\usepackage[a4paper,margin=25mm]{geometry}
\usepackage{amsmath,amssymb,mathtools}
\usepackage{xurl}
\usepackage[hidelinks]{hyperref}
\newcommand{\sref}[2]{\hyperlink{source:#1:#2}{[S#2]}}
\setlength{\emergencystretch}{3em}
\begin{document}
% problemId:30007551
\section*{Redirecting technological change}\label{30007551}
\subsection{Definitions and mathematical statement}
Fix two research directions \(j=1,2\), such as ordinary private productivity and a declared socially valuable health or environmental technology. Firms privately observe project quality \(q\), choose direction and research effort \(r\), and obtain innovation arrival rate \(\lambda_j(r,q)\). Private rewards \(v_j(q)\) differ from social rewards \(V_j(q)\), including measured spillovers and external damages. A policymaker observes imperfect signals \(s\), can offer subsidies \(t_j(s,r)\), and faces a stated budget and an explicit probability of political capture or biased project selection.
Determine a feasible subsidy or procurement rule maximizing expected discounted social surplus net of research resources, fiscal costs, and capture losses. Require firms’ participation and effort/direction incentive compatibility for every type in the declared quality support. Preserve a specified minimum experimentation probability for projects with uncertain social value, or incorporate the information value of experimentation directly in the objective.
Identify the private-response, spillover, learning, and capture parameters needed for policy comparison, and distinguish uncertainty about project returns from discretionary favoritism. Evaluate rules against untargeted support and no-support benchmarks under equal budgets. The target is direction-specific innovation policy with explicit informational constraints; it does not assert that a policymaker knows socially optimal technology ex ante or that one favored sector should receive support universally.

\par\textbf{Formulation and scope.} The source explicitly asks about directing innovation toward social needs while preserving experimentation. The two-direction mechanism and capture constraints are adapted.

\par\textbf{Open-question provenance.} An explicit question or research agenda was verified at the source locator. This does not by itself certify that every later version remains unresolved \sref{30007551}{1}.

\par\textbf{Historical priority.} This is the earliest explicit relevant statement verified in this audit, not a claim of original priority; older literature and earlier versions may contain antecedents.

\subsection{Short English statement}
Design innovation-direction incentives that maximize social surplus under private information, fiscal limits, capture risk, and experimentation requirements.

\subsection{Sources}
\begin{itemize}\raggedright
\item[\textnormal{[S1]}]\hypertarget{source:30007551:1}{} Ufuk Akcigit. 2026. Creative destruction in economic growth. Section 7, “Open questions,” item (7).
\url{https://onlinelibrary.wiley.com/doi/10.1111/sjoe.70037}.
DOI: 10.1111/sjoe.70037.
{\small\textit{Earliest verified question/agenda reference; historical priority qualified below. The author explicitly asks how policy can redirect innovation toward social priorities when private and social returns differ while preserving experimentation.}}
\end{itemize}
\end{document}
